\scanpage{39}
\begingroup\makeatletter\tagsleft@true\makeatother
\[
\left(\underset{\substack{\unsure{\rotatebox{-90}{$\in$}|\rotatebox{-90}{$\ni$}}\\\unsure{\omega_\alpha|w_\alpha}}}{x}
<
\underset{\substack{\rotatebox{-90}{$\in$}\\\omega_\beta}}{y}\right)
\Longleftrightarrow
\Biggl(\alpha\ne\beta,\quad
\underbrace{x\underset{\{\Pi_\alpha,\Pi_\beta\}}{<}y}_{
  \inserted{\substack{\text{i.e. }\deleted{\unsure{x}}\\\Sigma_\alpha^x\cap\Sigma_\beta^{\sigma y}=\varnothing}}}
\Biggr)
\tag*{(26)}
\]
% slanted
\margin{il faut vérifier la \unsure{transitivité} de cette relation -- c'est immédiat}
Alors $\sigma$ est un antiautomorphisme (involutif sans pt fixe) de $C$. L'application
\[
\begin{aligned}
C&\xrightarrow{\unsure{\varphi|\phi}}{\unsure{\mathfrak{P}|\mathcal{P}}}_{\unsure{\mathrm{finies}|\mathrm{ferm\acute{e}s}|\mathrm{foncier}}}(X)\\
(\alpha,x)&\longmapsto\Sigma_\alpha^x
\end{aligned}
\tag*{(27)}
\]
est une application croissante. Pour qu'\unsure{elles|elle} soit injective, et définisse un isomorphisme d'ens ordonnés de $C$ avec son image, il f et s. que $(\Pi_\alpha)_{\alpha\in A}$ soit strictement non \unsure{panaché|panachée}.

\unsure{Indépendamment} de $X$ :

\uline{Définition} On appelle \uline{\unsure{maquette|moquette} \deleted{\unsure{de}} \unsure{ordonnée}} (de \unsure{décompositions} \unsure{$n$-aires|$n$ aires}) un ens \deleted{\unsure{ordonné}} \inserted{fini} $C$, muni d'une involution \inserted{$\sigma$} sans pt fixe, et d'une relation d'ordre, ayant la \unsure{propriété|propriétés} \unsure{suivante|suivantes} :

\deleted{\unsure{a)} Si $x,y\in C$ et $x<y$, alors $x$ et $y$}
\deleted{Si $x,y\in C$, $\unsure{\sigma x\ne y|\sigma x=y}$, $x$ et $y$ non \unsure{congrus} sous $\sigma$, alors on n'a pas $x<y$ non \unsure{congrus} \unsure{mod} $C$,}
\deleted{\unsure{b)}} \deleted{\unsure{et}} \deleted{\unsure{b)}}
% slanted
\margin{NB implique que $x,y$ \unsure{compar.} $\Longrightarrow$ $x,y$ non \unsure{congrus} mod $\sigma$}
\[
\left\{\begin{aligned}
\text{a)}\;&\text{Si }x,y\in C,\quad\deleted{\unsure{x\ne y}},\quad
\text{alors }\deleted{\sigma x,y\text{ et }x,\sigma y}\ 
\inserted{x,y\text{ comparables}}\\
&\text{sss }\ \sigma x,y\ \text{non comparables}\\
\text{b)}\;&\text{Si }x,y\in C,\ \deleted{\text{comparables, et }\sigma y<\sigma x}\
\text{ alors }\ x<y\Longrightarrow\sigma y<\sigma x
\end{aligned}\right.
\tag*{(\unsure{28})}
\]
\uline{NB} Cela \unsure{indique|implique} que $\sigma$ est un antiautomorphisme de l'ens ordonné $C$

\endgroup
\clearpage
% archive sheet number: 38
